%% LyX 2.4.4 created this file.  For more info, see https://www.lyx.org/.
%% Do not edit unless you really know what you are doing.
\documentclass[english]{article}
\usepackage[T1]{fontenc}
\usepackage[latin9]{inputenc}
\usepackage{enumitem}
\usepackage{amsmath}
\usepackage{amsthm}
\usepackage{geometry}
\geometry{verbose,tmargin=1in,bmargin=1in,lmargin=1in,rmargin=1in}

\makeatletter
%%%%%%%%%%%%%%%%%%%%%%%%%%%%%% Textclass specific LaTeX commands.
\numberwithin{equation}{section}
\numberwithin{figure}{section}
\newlength{\lyxlabelwidth}      % auxiliary length 

%%%%%%%%%%%%%%%%%%%%%%%%%%%%%% User specified LaTeX commands.
\usepackage{fancyhdr}
\usepackage{lastpage}
\pagestyle{fancy}
\usepackage{enumitem}
\usepackage{ifthen}
%sets math fonts to be more readable
\everymath=\expandafter{\the\everymath\displaystyle}
%\DeclareMathSizes{10}{10}{10}{10}
\usepackage{multicol}
% Added by lyx2lyx
\setlength{\parskip}{\smallskipamount}
\setlength{\parindent}{0pt}

\makeatother

\usepackage{babel}
\begin{document}
\renewcommand{\author}{Dr.\,James Ashe}

\newcommand{\classNum}{232}

\newcommand{\classSec}{A and B}

\newcommand{\nameType}{}

\newcommand{\examType}{Exam 1 Review}

\renewcommand{\date}{Exam 1 will be be given on February $18^{\text{th}}$, 2012} %%\currenttime, \today

%%First page's header%%%%%%%%%%%%%%%%%%%%%%
\fancypagestyle{firststyle} %%header settings for first pagr
{
	\fancyhead[L]{MTH \classNum{}(\classSec{}) \author{}\\
	\textbf{\examType{}} \date}
	\fancyhead[C]{\textbf{\nameType}}
	\setlength{\headsep}{14pt}
}
%%%%%%%%%%%%%%%%%%%%%%%%%%%%%%%%%%%%%%%%%%

%%Next page's header%%%%%%%%%%%%%%%%%%%%%%%
\setlist{nolistsep}
\lhead{\classNum{}(\classSec{})} 
\chead{\textbf{\nameType}} 
\cfoot{\thepage{}~of~\pageref{LastPage}} 
\setlength{\headsep}{5pt} 
%%%%%%%%%%%%%%%%%%%%%%%%%%%%%%%%%%%%%%%%%%%

%%Various settings; see the preamble for other settings%%%%%
%%\setenumerate[0]{leftmargin=12pt,itemindent=0pt,labelwidth=0pt}
\renewcommand{\labelenumi}{\textbf{\arabic{enumi}.}}
\renewcommand{\labelenumii}{\textbf{\alph{enumii}.}}
\thispagestyle{firststyle}
%%%%%%%%%%%%%%%%%%%%%%%%%%%%%%%%%%%%%%%%%%

\vfill{}

\subsection*{4.8: Antiderivative}
\begin{itemize}
\item Indefinite integrals involving powers and basic trig formulas. \textbf{HW
1-8}
\item Initial value problems. \textbf{HW 9-15}
\end{itemize}

\subsection*{5.1: Approximating Areas Under the Curves}
\begin{itemize}
\item Calculating left and right Riemann sums. \textbf{HW 2, 5-7, 9}
\item Reading and writing sigma notation. \textbf{HW 10-11}
\end{itemize}

\subsection*{5.2: Definite Integrals. }
\begin{itemize}
\item Evaluating definite integrals using geometry. \textbf{HW 7-9, 17}
\item Properties of definite integrals

\begin{itemize}
\item Know the difference between finding the \emph{area }and \emph{net
area }of a region\emph{. }\textbf{HW 1, 19}
\item Evaluating definite integrals from graphs. \textbf{HW 10-12}
\item Evaluating definite integrals from given values. \textbf{HW 13-15}
\end{itemize}
\end{itemize}

\subsection*{5.3: Fundamental Theorem of Calculus}
\begin{itemize}
\item Applying the FTC, and understanding the inverse relationship between
differentiation and integration. \textbf{HW 1-5, 17-18}
\item Evaluating definite integrals for explicit values. \textbf{HW 7-12,
14, 19}
\item Find the area bounded by a curve. \textbf{HW 13, 20}

\begin{itemize}
\item Find the \emph{net area }of the region bounded by the curve and the
\emph{x-}axis on an interval, i.e., the definite integral on that
interval.
\item Find the \emph{area }of the regions bounded by the curve and the \emph{x-}axis
on an interval, i.e., calculate all area as positive.
\end{itemize}
\end{itemize}

\subsection*{5.4: Working with Integrals}
\begin{itemize}
\item Using properties of even and odd functions. \textbf{HW 1-2, 4-8, 19}
\item Finding the average value of a function. \textbf{HW 9-14, 18}
\end{itemize}

\subsection*{5.5: Substitution Rule}
\begin{itemize}
\item Using the substitution rule to find the indefinite integral. You \emph{will
not necessarily be told which problems require this rule. }\textbf{HW
1-10}
\end{itemize}

\subsection*{Some general notes.}
\begin{itemize}
\item You will be given the necessary formulas from section 4.8. 
\item The quiz and homework are good reviews. You can review pass due HW
assignments in MyMathLab (MML) by going to ``Gradebook'' and clicking
on ``Review''. Here you can practice more problems by clicking on
``Similar Exercise''.
\item More practice problems can be found in MML in ``Study Plan'' or
your textbook.
\item There will be 15-20 problems including some ``easy'' problems involving
properties of integrals which can worked quickly. 
\end{itemize}

\end{document}
